\documentclass[12pt,a4paper]{article}

\usepackage[utf8]{inputenc}
\usepackage[T1]{fontenc}
\usepackage[brazilian]{babel}
\usepackage[a4paper,margin=2.5cm]{geometry}
\usepackage{graphicx}
\usepackage{xcolor}
\usepackage{listings}
\usepackage{hyperref}
\usepackage{parskip}
\usepackage{fancyhdr}
\usepackage{array}
\usepackage{titlesec}

\definecolor{multivixred}{RGB}{166,27,37}
\definecolor{codebg}{RGB}{245,245,245}
\definecolor{codegray}{RGB}{110,110,110}
\definecolor{codeblue}{RGB}{0,90,160}

\titleformat{\section}{\large\bfseries\color{multivixred}}{\thesection}{1em}{}
\titleformat{\subsection}{\normalsize\bfseries}{\thesubsection}{1em}{}

\lstdefinelanguage{csharp}{
  morekeywords={class,public,private,protected,static,void,string,int,bool,
    new,base,this,using,namespace,return,if,else,foreach,in,async,await,Task,
    readonly,get,set},
  sensitive=true,
  morecomment=[l]{//},
  morestring=[b]"
}

\lstset{
  language=csharp,
  backgroundcolor=\color{codebg},
  basicstyle=\ttfamily\footnotesize,
  keywordstyle=\color{codeblue}\bfseries,
  commentstyle=\color{codegray}\itshape,
  stringstyle=\color{multivixred},
  numbers=left,
  numberstyle=\tiny\color{codegray},
  frame=single,
  rulecolor=\color{codegray},
  breaklines=true,
  showstringspaces=false,
  tabsize=2,
  extendedchars=true,
  literate=
    {á}{{\'a}}1 {à}{{\`a}}1 {ã}{{\~a}}1 {â}{{\^a}}1
    {é}{{\'e}}1 {ê}{{\^e}}1
    {í}{{\'i}}1
    {ó}{{\'o}}1 {õ}{{\~o}}1 {ô}{{\^o}}1
    {ú}{{\'u}}1 {ü}{{\"u}}1
    {ç}{{\c c}}1
    {Á}{{\'A}}1 {À}{{\`A}}1 {Ã}{{\~A}}1 {Â}{{\^A}}1
    {É}{{\'E}}1 {Ê}{{\^E}}1
    {Í}{{\'I}}1
    {Ó}{{\'O}}1 {Õ}{{\~O}}1 {Ô}{{\^O}}1
    {Ú}{{\'U}}1 {Ü}{{\"U}}1
    {Ç}{{\c C}}1
}

\pagestyle{fancy}
\fancyhf{}
\renewcommand{\headrulewidth}{0.4pt}
\fancyhead[L]{Programação Distribuída e Paralela}
\fancyhead[R]{Trabalho I}
\fancyfoot[C]{\thepage}

\hypersetup{
  colorlinks=true,
  linkcolor=multivixred,
  urlcolor=multivixred,
}

\begin{document}

% ---------- Capa ----------
\begin{titlepage}
  \centering
  \includegraphics[width=0.55\textwidth]{multivix-logo.png}\\[2.5cm]

  {\Large Faculdade Multivix -- Serra/ES}\\[0.3cm]
  {\large Curso de Engenharia da Computação}\\[2.5cm]

  {\Huge\bfseries Trabalho I}\\[0.4cm]
  {\Large Programação Distribuída e Paralela}\\[0.4cm]
  {\large Sistema de diálogo entre personagens de RPG com\\
   Filas, Tasks e métodos síncronos/assíncronos}\\[3cm]

  \begin{tabular}{ll}
    \textbf{Integrantes:} & Luiz Fernando Pereira Rojas dos Santos \\
                           & Tadeu Augusto Vilela da Silva \\[0.3cm]
    \textbf{Curso:}       & Engenharia da Computação \\
    \textbf{Turma:}       & 2026.2 \\
    \textbf{Professor:}   & Jacques Oliveira \\
    \textbf{Instituição:} & Multivix Serra -- Serra, ES \\
  \end{tabular}

  \vfill
  {\large Serra -- ES, \today}
\end{titlepage}

% ---------- Conteúdo ----------

\section{Objetivo do trabalho}

O trabalho solicitado pelo professor Jacques Oliveira consiste em implementar
um sistema simples de diálogo entre dois personagens no estilo RPG, de forma
caractere a caractere, utilizando os seguintes recursos da linguagem C\#:

\begin{itemize}
  \item \textbf{Filas} (\texttt{Queue<T>}) para armazenar as falas de cada personagem;
  \item \textbf{Métodos síncronos e assíncronos} para comparar os dois modelos de execução;
  \item \textbf{Tasks} (\texttt{Task}, \texttt{async}, \texttt{await}) para observar, na
    prática, o comportamento de operações concorrentes e paralelas e o controle de
    sincronização entre elas;
  \item \textbf{Orientação a Objetos} (classe, herança e instância), por se tratar de um
    exercício também de programação distribuída.
\end{itemize}

\section{Repositório do código}

O código-fonte completo está publicado, de forma pública, no GitHub:

\begin{center}
  \url{https://github.com/LuizRojas/trabalho-programacao-paralela}
\end{center}

O projeto é uma aplicação console em \textbf{.NET} e pode ser executada com:

\begin{lstlisting}[language=bash,numbers=none]
dotnet run
\end{lstlisting}

\section{O que foi implementado}

\subsection{Modelagem orientada a objetos}

Foi criada uma classe base \texttt{Personagem}, responsável por tudo o que é comum
a qualquer personagem do diálogo: nome, tipo, a fila de falas e os métodos de
``falar'' (síncrono e assíncrono). As classes \texttt{Mago} e \texttt{Elfo} \textbf{herdam}
de \texttt{Personagem} e, no construtor, apenas enfileiram as falas específicas de
cada uma. Cada personagem usado no programa (\texttt{Mago("Gandalf")} e
\texttt{Elfo("Link")}) é uma \textbf{instância} independente dessas classes, com sua
própria fila de mensagens.

\begin{lstlisting}
public class Personagem
{
    public string nome { get; }
    public string tipo { get; }
    protected Queue<string> mensagens;
    ...
}

public class Mago : Personagem { ... }  // herança
public class Elfo : Personagem { ... }  // herança
\end{lstlisting}

\subsection{Filas (Queue<T>)}

Cada personagem guarda suas falas em uma \texttt{Queue<string>}. A fila garante que
as mensagens sejam exibidas sempre na ordem em que foram cadastradas
(\texttt{FIFO -- First In, First Out}), que é exatamente o comportamento esperado
em um diálogo: a primeira fala cadastrada é a primeira a ser ``dita''.

\subsection{Método síncrono}

O método \texttt{DialogoSync} retira uma fala da fila e a imprime letra a letra de
forma \textbf{bloqueante}, usando \texttt{Thread.Sleep}. Enquanto uma fala está sendo
exibida, a aplicação não consegue fazer mais nada -- é a thread principal que fica
ocupada do início ao fim da fala.

\subsection{Método assíncrono (async/await/Task)}

O método \texttt{DialogoAsync} tem a mesma ideia, mas retorna \texttt{Task} e usa
\texttt{await Task.Delay(...)} em vez de \texttt{Thread.Sleep}. Isso faz com que a
espera entre cada caractere \textbf{não bloqueie} a thread: o controle é devolvido
ao chamador, que pode iniciar outras tasks em paralelo enquanto a atual está
``aguardando''.

\begin{lstlisting}
public async Task DialogoAsync(int tempoFala = 80)
{
    ...
    await _semaphore.WaitAsync();
    try
    {
        foreach (char c in msn)
        {
            Console.Write($"{c}");
            await Task.Delay(tempoFala);
        }
    }
    finally
    {
        _semaphore.Release();
    }
}
\end{lstlisting}

\subsection{Sincronização com SemaphoreSlim}

No \texttt{Program.cs}, várias chamadas a \texttt{DialogoAsync} são disparadas quase ao
mesmo tempo e aguardadas com \texttt{Task.WhenAll}, caracterizando execução
\textbf{concorrente} (várias tasks ``em andamento'' ao mesmo tempo). Se nenhum
controle fosse feito, os caracteres de falas diferentes poderiam se misturar no
console (ex.: \texttt{M}, de ``Mago'', intercalado com \texttt{E}, de ``Elfo'').

Para evitar isso, foi usado um \texttt{SemaphoreSlim(1, 1)}, que funciona como um
cadeado: apenas uma task por vez pode estar dentro da seção que imprime uma fala
completa. As demais tasks ficam aguardando (\texttt{await \_semaphore.WaitAsync()})
até a atual liberar o semáforo (\texttt{\_semaphore.Release()}), garantindo que cada
fala apareça inteira e organizada, mesmo com várias tasks rodando ao mesmo tempo.

\subsection{Disparo em paralelo com Task[] e Task.WhenAll}

\begin{lstlisting}
Task[] tarefas = {
    mago.DialogoAsync(600),
    elfo.DialogoAsync(),
    mago.DialogoAsync(1000),
    elfo.DialogoAsync(),
    mago.DialogoAsync(800),
    elfo.DialogoAsync()
};

await Task.WhenAll(tarefas);
\end{lstlisting}

Todas as tasks são criadas de uma vez (array \texttt{Task[]}) e o programa só
continua quando todas terminarem (\texttt{Task.WhenAll}), demonstrando o paralelismo
de forma controlada, em contraste com a execução puramente sequencial do bloco
síncrono anterior.

\section{Entendimento do grupo}

\subsection{Concorrência vs. paralelismo}

\textbf{Concorrência} é a capacidade de lidar com várias tarefas ``ao mesmo tempo''
do ponto de vista lógico (elas se alternam, podem estar pausadas aguardando
algo). \textbf{Paralelismo} é quando essas tarefas realmente executam
simultaneamente (em múltiplos núcleos). No trabalho, o uso de \texttt{async/await}
cria concorrência: as tasks de diálogo ficam alternando entre executar um
pedaço e aguardar (\texttt{Task.Delay}), sem travar a thread principal -- e é o
\texttt{Task.WhenAll} que efetivamente as executa de forma paralela/concorrente.

\subsection{Por que síncrono bloqueia e assíncrono não}

No método síncrono, \texttt{Thread.Sleep} realmente ``congela'' a thread até o
tempo passar. No método assíncrono, \texttt{await Task.Delay} devolve o controle
para quem chamou a task, permitindo que outras tasks avancem enquanto a atual
está esperando -- por isso é possível disparar os diálogos do \texttt{Mago} e do
\texttt{Elfo} juntos, sem que um precise esperar o outro terminar por completo.

\subsection{Por que foi necessário um semáforo}

Mesmo com tasks assíncronas, o \texttt{Console} é um recurso compartilhado. Sem
sincronização, duas tasks poderiam escrever no console ao mesmo tempo e misturar
os caracteres das falas. O \texttt{SemaphoreSlim} resolve isso garantindo exclusão
mútua durante a escrita de cada fala, sem bloquear a thread (porque
\texttt{WaitAsync} também é assíncrono).

\section{Resultado obtido}

A saída abaixo foi gerada executando \texttt{dotnet run} e corresponde ao
resultado esperado no enunciado:

\begin{lstlisting}[numbers=none,basicstyle=\ttfamily\scriptsize]
Mago: Quem é você viajante ?...
Elfo: Sou link do jogo zelda...
Mago: OK, sou Gandalf do senhor dos anéis...
Mago: Se diz ser quem é, então diga as palavras secretas comigo...
Mago: E as portas se abrirão...
Mago: Anah  Elfo: Milac  Mago: Passi-cii-fiik  Elfo: Teish  Mago: Perim  Elfo: Salaa-saa-aah

***=================================***
Portas abertas!
\end{lstlisting}

\section{Conclusão}

O trabalho permitiu observar, na prática, a diferença entre execução síncrona
(bloqueante) e assíncrona (não bloqueante), além do uso de \texttt{Task} para
paralelizar o diálogo de dois personagens que, mesmo ``conversando ao mesmo
tempo'', precisam de sincronização (semáforo) para que a saída no console
permaneça legível e organizada -- exatamente o comportamento clássico de
concorrência sobre um recurso compartilhado.

\end{document}
